\IfFileExists{stacks-project.cls}{%
\documentclass{stacks-project}
}{%
\documentclass{amsart}
}

% For dealing with references we use the comment environment
\usepackage{verbatim}
\newenvironment{reference}{\comment}{\endcomment}
%\newenvironment{reference}{}{}
\newenvironment{slogan}{\comment}{\endcomment}
\newenvironment{history}{\comment}{\endcomment}

% For commutative diagrams we use Xy-pic
\usepackage[all]{xy}

% We use 2cell for 2-commutative diagrams.
\xyoption{2cell}
\UseAllTwocells

% We use multicol for the list of chapters between chapters
\usepackage{multicol}

% This is generally recommended for better output
\usepackage{lmodern}
\usepackage[T1]{fontenc}

% For cross-file-references
\usepackage{xr-hyper}

% Package for hypertext links:
\usepackage{hyperref}

% For any local file, say "hello.tex" you want to link to please
% use \externaldocument[hello-]{hello}
\externaldocument[introduction-]{introduction}
\externaldocument[conventions-]{conventions}
\externaldocument[sets-]{sets}
\externaldocument[categories-]{categories}
\externaldocument[topology-]{topology}
\externaldocument[sheaves-]{sheaves}
\externaldocument[sites-]{sites}
\externaldocument[stacks-]{stacks}
\externaldocument[fields-]{fields}
\externaldocument[algebra-]{algebra}
\externaldocument[brauer-]{brauer}
\externaldocument[homology-]{homology}
\externaldocument[derived-]{derived}
\externaldocument[simplicial-]{simplicial}
\externaldocument[more-algebra-]{more-algebra}
\externaldocument[smoothing-]{smoothing}
\externaldocument[modules-]{modules}
\externaldocument[sites-modules-]{sites-modules}
\externaldocument[injectives-]{injectives}
\externaldocument[cohomology-]{cohomology}
\externaldocument[sites-cohomology-]{sites-cohomology}
\externaldocument[dga-]{dga}
\externaldocument[dpa-]{dpa}
\externaldocument[sdga-]{sdga}
\externaldocument[hypercovering-]{hypercovering}
\externaldocument[schemes-]{schemes}
\externaldocument[constructions-]{constructions}
\externaldocument[properties-]{properties}
\externaldocument[morphisms-]{morphisms}
\externaldocument[coherent-]{coherent}
\externaldocument[divisors-]{divisors}
\externaldocument[limits-]{limits}
\externaldocument[varieties-]{varieties}
\externaldocument[topologies-]{topologies}
\externaldocument[descent-]{descent}
\externaldocument[perfect-]{perfect}
\externaldocument[more-morphisms-]{more-morphisms}
\externaldocument[flat-]{flat}
\externaldocument[groupoids-]{groupoids}
\externaldocument[more-groupoids-]{more-groupoids}
\externaldocument[etale-]{etale}
\externaldocument[chow-]{chow}
\externaldocument[intersection-]{intersection}
\externaldocument[pic-]{pic}
\externaldocument[weil-]{weil}
\externaldocument[adequate-]{adequate}
\externaldocument[dualizing-]{dualizing}
\externaldocument[duality-]{duality}
\externaldocument[discriminant-]{discriminant}
\externaldocument[derham-]{derham}
\externaldocument[local-cohomology-]{local-cohomology}
\externaldocument[algebraization-]{algebraization}
\externaldocument[curves-]{curves}
\externaldocument[resolve-]{resolve}
\externaldocument[models-]{models}
\externaldocument[functors-]{functors}
\externaldocument[equiv-]{equiv}
\externaldocument[pione-]{pione}
\externaldocument[etale-cohomology-]{etale-cohomology}
\externaldocument[proetale-]{proetale}
\externaldocument[relative-cycles-]{relative-cycles}
\externaldocument[more-etale-]{more-etale}
\externaldocument[trace-]{trace}
\externaldocument[crystalline-]{crystalline}
\externaldocument[spaces-]{spaces}
\externaldocument[spaces-properties-]{spaces-properties}
\externaldocument[spaces-morphisms-]{spaces-morphisms}
\externaldocument[decent-spaces-]{decent-spaces}
\externaldocument[spaces-cohomology-]{spaces-cohomology}
\externaldocument[spaces-limits-]{spaces-limits}
\externaldocument[spaces-divisors-]{spaces-divisors}
\externaldocument[spaces-over-fields-]{spaces-over-fields}
\externaldocument[spaces-topologies-]{spaces-topologies}
\externaldocument[spaces-descent-]{spaces-descent}
\externaldocument[spaces-perfect-]{spaces-perfect}
\externaldocument[spaces-more-morphisms-]{spaces-more-morphisms}
\externaldocument[spaces-flat-]{spaces-flat}
\externaldocument[spaces-groupoids-]{spaces-groupoids}
\externaldocument[spaces-more-groupoids-]{spaces-more-groupoids}
\externaldocument[bootstrap-]{bootstrap}
\externaldocument[spaces-pushouts-]{spaces-pushouts}
\externaldocument[spaces-chow-]{spaces-chow}
\externaldocument[groupoids-quotients-]{groupoids-quotients}
\externaldocument[spaces-more-cohomology-]{spaces-more-cohomology}
\externaldocument[spaces-simplicial-]{spaces-simplicial}
\externaldocument[spaces-duality-]{spaces-duality}
\externaldocument[formal-spaces-]{formal-spaces}
\externaldocument[restricted-]{restricted}
\externaldocument[spaces-resolve-]{spaces-resolve}
\externaldocument[formal-defos-]{formal-defos}
\externaldocument[defos-]{defos}
\externaldocument[cotangent-]{cotangent}
\externaldocument[examples-defos-]{examples-defos}
\externaldocument[algebraic-]{algebraic}
\externaldocument[examples-stacks-]{examples-stacks}
\externaldocument[stacks-sheaves-]{stacks-sheaves}
\externaldocument[criteria-]{criteria}
\externaldocument[artin-]{artin}
\externaldocument[quot-]{quot}
\externaldocument[stacks-properties-]{stacks-properties}
\externaldocument[stacks-morphisms-]{stacks-morphisms}
\externaldocument[stacks-limits-]{stacks-limits}
\externaldocument[stacks-cohomology-]{stacks-cohomology}
\externaldocument[stacks-perfect-]{stacks-perfect}
\externaldocument[stacks-introduction-]{stacks-introduction}
\externaldocument[stacks-more-morphisms-]{stacks-more-morphisms}
\externaldocument[stacks-geometry-]{stacks-geometry}
\externaldocument[moduli-]{moduli}
\externaldocument[moduli-curves-]{moduli-curves}
\externaldocument[examples-]{examples}
\externaldocument[exercises-]{exercises}
\externaldocument[guide-]{guide}
\externaldocument[desirables-]{desirables}
\externaldocument[coding-]{coding}
\externaldocument[obsolete-]{obsolete}
\externaldocument[fdl-]{fdl}
\externaldocument[index-]{index}

% Theorem environments.
%
\theoremstyle{plain}
\newtheorem{theorem}[subsection]{Theorem}
\newtheorem{proposition}[subsection]{Proposition}
\newtheorem{lemma}[subsection]{Lemma}

\theoremstyle{definition}
\newtheorem{definition}[subsection]{Definition}
\newtheorem{example}[subsection]{Example}
\newtheorem{exercise}[subsection]{Exercise}
\newtheorem{situation}[subsection]{Situation}

\theoremstyle{remark}
\newtheorem{remark}[subsection]{Remark}
\newtheorem{remarks}[subsection]{Remarks}

\numberwithin{equation}{subsection}

% Macros
%
\def\lim{\mathop{\mathrm{lim}}\nolimits}
\def\colim{\mathop{\mathrm{colim}}\nolimits}
\def\Spec{\mathop{\mathrm{Spec}}}
\def\Hom{\mathop{\mathrm{Hom}}\nolimits}
\def\Ext{\mathop{\mathrm{Ext}}\nolimits}
\def\SheafHom{\mathop{\mathcal{H}\!\mathit{om}}\nolimits}
\def\SheafExt{\mathop{\mathcal{E}\!\mathit{xt}}\nolimits}
\def\Sch{\mathit{Sch}}
\def\Mor{\mathop{\mathrm{Mor}}\nolimits}
\def\Ob{\mathop{\mathrm{Ob}}\nolimits}
\def\Sh{\mathop{\mathit{Sh}}\nolimits}
\def\NL{\mathop{N\!L}\nolimits}
\def\CH{\mathop{\mathrm{CH}}\nolimits}
\def\proetale{{pro\text{-}\acute{e}tale}}
\def\etale{{\acute{e}tale}}
\def\QCoh{\mathit{QCoh}}
\def\Ker{\mathop{\mathrm{Ker}}}
\def\Im{\mathop{\mathrm{Im}}}
\def\Coker{\mathop{\mathrm{Coker}}}
\def\Coim{\mathop{\mathrm{Coim}}}

% Boxtimes
%
\DeclareMathSymbol{\boxtimes}{\mathbin}{AMSa}{"02}

%
% Macros for moduli stacks/spaces
%
\def\QCohstack{\mathcal{QC}\!\mathit{oh}}
\def\Cohstack{\mathcal{C}\!\mathit{oh}}
\def\Spacesstack{\mathcal{S}\!\mathit{paces}}
\def\Quotfunctor{\mathrm{Quot}}
\def\Hilbfunctor{\mathrm{Hilb}}
\def\Curvesstack{\mathcal{C}\!\mathit{urves}}
\def\Polarizedstack{\mathcal{P}\!\mathit{olarized}}
\def\Complexesstack{\mathcal{C}\!\mathit{omplexes}}
% \Pic is the operator that assigns to X its picard group, usage \Pic(X)
% \Picardstack_{X/B} denotes the Picard stack of X over B
% \Picardfunctor_{X/B} denotes the Picard functor of X over B
\def\Pic{\mathop{\mathrm{Pic}}\nolimits}
\def\Picardstack{\mathcal{P}\!\mathit{ic}}
\def\Picardfunctor{\mathrm{Pic}}
\def\Deformationcategory{\mathcal{D}\!\mathit{ef}}
\usepackage{cleveref}
\crefname{originalreference}{original reference}{original references}
\Crefname{originalreference}{Original reference}{Original references}
\makeatletter
\newlabel{lemma-compare-simplicial-objects}{{0DAF}{source}{}{sourcepointer-1}{}}
\newlabel{lemma-compare-simplicial-objects@cref}{{[originalreference][0][]0DAF}{[1][0][]source}{}{}{}}
\makeatother
\crefname{originalreference}{original reference}{original references}
\Crefname{originalreference}{Original reference}{Original references}
\makeatletter
\newlabel{lemma-compare-simplicial-objects-fppf-etale}{{0DH5}{source}{}{sourcepointer-2}{}}
\newlabel{lemma-compare-simplicial-objects-fppf-etale@cref}{{[originalreference][0][]0DH5}{[1][0][]source}{}{}{}}
\makeatother
\crefname{originalreference}{original reference}{original references}
\Crefname{originalreference}{Original reference}{Original references}
\makeatletter
\newlabel{lemma-direct-image-morphism-simplicial-sites-modules}{{0DH1}{source}{}{sourcepointer-3}{}}
\newlabel{lemma-direct-image-morphism-simplicial-sites-modules@cref}{{[originalreference][0][]0DH1}{[1][0][]source}{}{}{}}
\makeatother
\crefname{originalreference}{original reference}{original references}
\Crefname{originalreference}{Original reference}{Original references}
\makeatletter
\newlabel{lemma-hypercovering-X-simple-descent-bounded-modules}{{0DAC}{source}{}{sourcepointer-4}{}}
\newlabel{lemma-hypercovering-X-simple-descent-bounded-modules@cref}{{[originalreference][0][]0DAC}{[1][0][]source}{}{}{}}
\makeatother
\crefname{originalreference}{original reference}{original references}
\Crefname{originalreference}{Original reference}{Original references}
\makeatletter
\newlabel{remark-augmentation-ringed}{{0D9W}{source}{}{sourcepointer-5}{}}
\newlabel{remark-augmentation-ringed@cref}{{[originalreference][0][]0D9W}{[1][0][]source}{}{}{}}
\makeatother
\crefname{originalsection}{original section}{original sections}
\Crefname{originalsection}{Original section}{Original sections}
\makeatletter
\newlabel{section-simplicial-algebraic-spaces}{{0DE7}{source}{}{sourcepointer-6}{}}
\newlabel{section-simplicial-algebraic-spaces@cref}{{[originalsection][0][]0DE7}{[1][0][]source}{}{}{}}
\makeatother
\crefname{originallemma}{original lemma}{original lemmas}
\Crefname{originallemma}{Original lemma}{Original lemmas}
\makeatletter
\newlabel{spaces-more-cohomology-lemma-cohomological-descent-etale-fppf-modules}{{0DGP}{source}{}{sourcepointer-7}{}}
\newlabel{spaces-more-cohomology-lemma-cohomological-descent-etale-fppf-modules@cref}{{[originallemma][0][]0DGP}{[1][0][]source}{}{}{}}
\makeatother
\title{Cohomological descent for fppf hypercoverings on quasi-coherent modules: F -> Ra\_*a\textasciicircum{}*F is an isomorphism}
\author{The Stacks Project Authors}
\date{stacks-project:spaces-simplicial.tex:local-snapshot-20261012; GNU FDL 1.2 or later}
\begin{document}
\maketitle
\noindent\textit{Editorial note.} All delivered mathematical text is the author's own, copied from the pinned Stacks Project chapter file sources/stacks-spaces-simplicial/pinned-original-source.tex; no statement or proof step was rewritten, reordered or removed.\par\medskip
\noindent\textit{Editorial note.} The counted proof is the author's complete argument: the flatness identification La\_\{fppf\}\textasciicircum{}* = a\_\{fppf\}\textasciicircum{}* is justified by a\_\{fppf\}\textasciicircum{}\{-1\}O\_big = O\_big,total (source lines 7660-7664, referring to remark-augmentation-ringed); this lemma upgrades card L46172 (equivalence on QCoh) to cohomological descent for the derived pushforward, over the diagram established on card L46173; no verification is omitted.\par\medskip
\noindent\textit{Editorial note.} The label lemma-cohomological-descent-qcoh-for-fppf-hypercovering names the counted statement, and the proof environment immediately following it in the source is delivered verbatim as its proof.\par\medskip
\noindent\textit{Editorial note.} The counted statement is carried by the author's own lemma environment in the chapter source.\par\medskip
\noindent\textit{Allowed standard prerequisites (editorial).}\begin{itemize}
\item The fppf-etale comparison diagram of topoi. Same chapter, \textbackslash{}label\{lemma-compare-simplicial-objects-fppf-etale\} (card L46173). Proof lines 7631-7632 and 7650-7652: the commutative diagram with h, a\_\{fppf\}, h\_\{-1\}, a used throughout the computation.
\item Componentwise cohomological descent for QCoh modules. More on Cohomology of Spaces, \textbackslash{}label\{spaces-more-cohomology-lemma-cohomological-descent-etale-fppf-modules\}. Proof lines 7635-7638 and 7667-7673: F\_n = a\_n\textasciicircum{}*F satisfies F\_n = Rh\_\{n,*\}h\_n\textasciicircum{}*F\_n, and the closing equality Rh\_\{-1,*\}(h\_\{-1\})\textasciicircum{}*F = F.
\item Direct image on the total site is componentwise (modules). Same chapter, \textbackslash{}label\{lemma-direct-image-morphism-simplicial-sites-modules\}. Proof lines 7639-7640: a\textasciicircum{}*F = Rh\_*h\textasciicircum{}*a\textasciicircum{}*F obtained by assembling the component identities.
\item Simple descent for bounded modules along fppf hypercoverings. Same chapter, \textbackslash{}label\{lemma-hypercovering-X-simple-descent-bounded-modules\}. Proof lines 7655-7660: Ra\_\{fppf,*\}a\_\{fppf\}\textasciicircum{}* = id with La\_\{fppf\}\textasciicircum{}* = a\_\{fppf\}\textasciicircum{}* since a\_\{fppf\} is flat (a\_\{fppf\}\textasciicircum{}\{-1\}O\_big = O\_big,total, remark-augmentation-ringed).
\end{itemize}\medskip
\ifdefined\hypertarget\hypertarget{sourcepointer-1}{}\fi
\noindent\textit{Source reference.} Original reference 0DAF: sources/stacks-spaces-simplicial/pinned-original-source.tex, source lines 5910--5910. Prerequisite pointer invoked at proof line(s) 7664; it names a standard fact of the same chapter that the counted argument uses, and every step of the counted proof built on it is retained verbatim in the delivered body.\par\medskip
\ifdefined\hypertarget\hypertarget{sourcepointer-2}{}\fi
\noindent\textit{Source reference.} Original reference 0DH5: sources/stacks-spaces-simplicial/pinned-original-source.tex, source lines 7233--7233. Prerequisite pointer invoked at proof line(s) 7642; it names a standard fact of the same chapter that the counted argument uses, and every step of the counted proof built on it is retained verbatim in the delivered body.\par\medskip
\ifdefined\hypertarget\hypertarget{sourcepointer-3}{}\fi
\noindent\textit{Source reference.} Original reference 0DH1: sources/stacks-spaces-simplicial/pinned-original-source.tex, source lines 1656--1656. Prerequisite pointer invoked at proof line(s) 7649; it names a standard fact of the same chapter that the counted argument uses, and every step of the counted proof built on it is retained verbatim in the delivered body.\par\medskip
\ifdefined\hypertarget\hypertarget{sourcepointer-4}{}\fi
\noindent\textit{Source reference.} Original reference 0DAC: sources/stacks-spaces-simplicial/pinned-original-source.tex, source lines 5226--5226. Prerequisite pointer invoked at proof line(s) 7665; it names a standard fact of the same chapter that the counted argument uses, and every step of the counted proof built on it is retained verbatim in the delivered body.\par\medskip
\ifdefined\hypertarget\hypertarget{sourcepointer-5}{}\fi
\noindent\textit{Source reference.} Original reference 0D9W: sources/stacks-spaces-simplicial/pinned-original-source.tex, source lines 4073--4073. Prerequisite pointer invoked at proof line(s) 7669; it names a standard fact of the same chapter that the counted argument uses, and every step of the counted proof built on it is retained verbatim in the delivered body.\par\medskip
\ifdefined\hypertarget\hypertarget{sourcepointer-6}{}\fi
\noindent\textit{Source reference.} Original section 0DE7: sources/stacks-spaces-simplicial/pinned-original-source.tex, source lines 6985--6985. Prerequisite pointer invoked at proof line(s) 7638; it names a standard fact of the same chapter that the counted argument uses, and every step of the counted proof built on it is retained verbatim in the delivered body.\par\medskip
\ifdefined\hypertarget\hypertarget{sourcepointer-7}{}\fi
\noindent\textit{Source reference.} Original lemma 0DGP: sources/stacks-chow/tags.txt, source lines 16879--16879. Prerequisite pointer invoked at proof line(s) 7647; it names a standard fact of a sibling chapter that the counted argument uses, and every step of the counted proof built on it is retained verbatim in the delivered body.\par\medskip
\begin{lemma}
\label{lemma-cohomological-descent-qcoh-for-fppf-hypercovering}
Let $S$ be a scheme. Let $X$ be an algebraic space over $S$.
Let $U$ be a simplicial algebraic space over $S$. Let $a : U \to X$
be an augmentation. If $a : U \to X$ is an fppf hypercovering of $X$,
then for $\mathcal{F}$ a quasi-coherent $\mathcal{O}_X$-module
the map
$$
\mathcal{F} \to Ra_*(a^*\mathcal{F})
$$
is an isomorphism. Here $a : \Sh(U_\etale) \to \Sh(X_\etale)$
is as in Section \ref{section-simplicial-algebraic-spaces}.
\end{lemma}

\begin{proof}
Consider the diagram of Lemma \ref{lemma-compare-simplicial-objects-fppf-etale}.
Let $\mathcal{F}_n = a_n^*\mathcal{F}$ be the $n$th component of
$a^*\mathcal{F}$. This is a quasi-coherent $\mathcal{O}_{U_n}$-module.
Then $\mathcal{F}_n = Rh_{n, *}h_n^*\mathcal{F}_n$
by More on Cohomology of Spaces, Lemma
\ref{spaces-more-cohomology-lemma-cohomological-descent-etale-fppf-modules}.
Hence $a^*\mathcal{F} = Rh_*h^*a^*\mathcal{F}$ by
Lemma \ref{lemma-direct-image-morphism-simplicial-sites-modules}.
We have
\begin{align*}
Ra_*(a^*\mathcal{F})
& =
Ra_*Rh_*h^*a^*\mathcal{F} \\
& =
Rh_{-1, *}Ra_{fppf, *}a_{fppf}^*(h_{-1})^*\mathcal{F} \\
& =
Rh_{-1, *}(h_{-1})^*\mathcal{F} \\
& =
\mathcal{F}
\end{align*}
The first equality by the discussion above, the second equality
because of the commutativity of the diagram in
Lemma \ref{lemma-compare-simplicial-objects}, the third equality by
Lemma \ref{lemma-hypercovering-X-simple-descent-bounded-modules}
as $U$ is a hypercovering of $X$ in $(\textit{Spaces}/S)_{fppf}$
and $La_{fppf}^* = a_{fppf}^*$ as $a_{fppf}$ is flat
(namely $a_{fppf}^{-1}\mathcal{O}_{big} = \mathcal{O}_{big, total}$,
see Remark \ref{remark-augmentation-ringed}), and
the last equality by the already used
More on Cohomology of Spaces, Lemma
\ref{spaces-more-cohomology-lemma-cohomological-descent-etale-fppf-modules}.
\end{proof}

\end{document}
